\documentclass[12pt]{article}
\usepackage[margin=0.75in]{geometry}
\usepackage{fontspec}
\setmainfont{Latin Modern Roman}
\usepackage{microtype}
\usepackage{fancyhdr}
\usepackage{titlesec}
\usepackage{enumitem}
\usepackage{xcolor}
\usepackage[hidelinks]{hyperref}
\setlength{\parindent}{1.1em}
\setlength{\parskip}{0.35em}
\setlength{\headheight}{15pt}
\titleformat{\section}{\large\bfseries\scshape}{}{0pt}{}
\titleformat{\subsection}{\normalsize\bfseries}{}{0pt}{}
\pagestyle{fancy}
\fancyhf{}
\fancyhead[L]{Whately Elements of Logic}
\fancyhead[R]{Teacher Guide}
\fancyfoot[C]{\thepage}
\renewcommand{\headrulewidth}{0.4pt}

\begin{document}
\begin{center}
{\LARGE\bfseries Teacher Guide: Week 6}\par\vspace{0.25em}
{\large\scshape Good Form, Bad Facts}\par\vspace{0.4em}
{Validity, Truth, and Soundness in Whately and \textit{Julius Caesar}}\par\vspace{0.6em}
\rule{0.82\textwidth}{0.4pt}
\end{center}

\section*{Week 6 Overview}
Week 6 corrects a possible misunderstanding from Week 5. Students have learned to test form through mood and figure. They now learn that valid form is not the same as truth. A valid argument can fail if one premise is false or unproved. A true conclusion can also be badly proved. The goal is a two-column habit of diagnosis: form test and fact test.

The lit-anchor remains with \textit{Julius Caesar}. Brutus's argument can be reconstructed in valid form, but Antony attacks the factual premise that Caesar's ambition would enslave Rome. This lets students see the distinction between formal validity and soundness inside a public, literary, and moral argument.

\section*{Essential Question}
Can an argument be logically valid and still fail to prove what we should believe?

\section*{Student-Facing Materials}
\begin{itemize}[leftmargin=*]
\item \textit{Good Form, Bad Facts} -- student reading handout.
\item \textit{Week 6 Argument Lab: Good Form, Bad Facts} -- student lab for form-test/fact-test diagnosis.
\item \textit{Julius Caesar Lit-Anchor: Valid Form, Unproved Premise} -- Shakespeare anchor passage for testing Brutus's premise and Antony's answer.
\end{itemize}

\section*{Learning Goals}
Students will be able to:
\begin{enumerate}[leftmargin=*]
\item Distinguish validity from truth.
\item Define soundness as valid form plus true premises.
\item Explain why a valid argument may still be unsound.
\item Explain why a true conclusion may still be badly proved.
\item Identify the premise pressure point in an argument.
\item Classify a premise as true, false, doubtful, or not yet established.
\item Diagnose arguments using a form-test/fact-test chart.
\item Apply the distinction to Brutus's and Antony's funeral speeches in \textit{Julius Caesar}.
\end{enumerate}

\section*{Key Vocabulary}
valid, invalid, true premise, false premise, doubtful premise, sound argument, unsound argument, form test, fact test, material truth, overclaim, premise pressure point

\section*{Core Teaching Emphasis}
Do not let students use \textit{valid} as a synonym for \textit{good}, \textit{true}, or \textit{persuasive}. Validity is about whether the conclusion follows. Soundness requires validity plus true premises. This week should form a durable habit for essays: after arranging an argument clearly, students must ask whether the premises are actually supported.

A useful classroom refrain: \textit{Form first, facts second, final judgment last.}

\section*{Weekly Lesson Sequence}

\subsection*{Day 1: Validity Is Not Truth}
\begin{itemize}[leftmargin=*]
\item Review Week 5's form test: Does the conclusion follow?
\item Introduce the fact test: Are the premises true or adequately supported?
\item Use the deliberately false bird/eagle example from the reading.
\item Ask students to explain why the argument is valid but unsound.
\end{itemize}

\subsection*{Day 2: False Premise, Doubtful Premise, Bad Proof}
\begin{itemize}[leftmargin=*]
\item Work through the difference between false premise and doubtful premise.
\item Use the ``true conclusion, bad proof'' example to show that being right is not the same as proving well.
\item Students complete Part III of the Argument Lab.
\item Require each diagnosis to name the specific failure instead of merely saying ``bad argument.''
\end{itemize}

\subsection*{Day 3: Two-Column Diagnosis}
\begin{itemize}[leftmargin=*]
\item Model the Form Test / Fact Test chart on the board.
\item Students choose one lab argument and complete Part IV.
\item Pair-share: one student explains form; the other challenges the fact test.
\item Close by asking what evidence would settle the pressure-point premise.
\end{itemize}

\subsection*{Day 4: Julius Caesar Lit-Anchor}
\begin{itemize}[leftmargin=*]
\item Read Brutus's Act 3, Scene 2 passage.
\item Reconstruct Brutus's argument in valid form.
\item Identify the pressure point: Caesar was a man whose ambition would enslave Rome.
\item Read Antony's evidence: captives/ransoms, tears for the poor, refusal of the crown.
\item Students begin the \textit{Julius Caesar Form Test / Fact Test Diagnosis} portfolio artifact.
\end{itemize}

\subsection*{Day 5: Soundness and Synthesis}
\begin{itemize}[leftmargin=*]
\item Students finish the lab or lit-anchor portfolio artifact.
\item Seminar question: Does Antony prove Brutus wrong, or does he make Brutus's premise doubtful?
\item Connect to writing: a good essay needs not only a clean structure but also supported premises.
\item Exit writing: ``An argument can be valid but not sound when...'' Require a Caesar example.
\end{itemize}

\section*{Answer Guidance: Brutus}
Recommended diagnosis:
\begin{itemize}[leftmargin=*]
\item \textbf{Argument:} All men whose ambition would enslave Rome deserve death for Rome's freedom; Caesar was such a man; therefore Caesar deserved death for Rome's freedom.
\item \textbf{Form test:} Valid, if reconstructed as All M are P; All S are M; therefore All S are P.
\item \textbf{Fact test:} The minor premise is doubtful or not yet established. Brutus asserts Caesar's dangerous ambition but does not fully prove it in the funeral speech.
\item \textbf{Final diagnosis:} Not yet proved, or unsound if students argue the minor premise is false. Reward careful distinction.
\end{itemize}

\section*{Answer Guidance: Antony}
Antony attacks Brutus's pressure-point premise. His main evidence:
\begin{itemize}[leftmargin=*]
\item Caesar brought captives and ransom money into Rome's public coffers.
\item Caesar wept for the poor.
\item Caesar refused the crown three times at the Lupercal.
\end{itemize}

This evidence does not automatically prove Caesar had no ambition at all. It does make Brutus's specific premise doubtful: that Caesar's ambition was tyrannical and enslaving. Strong students should see that Antony weakens Brutus's soundness claim even if he does not produce a complete formal proof of Caesar's innocence.

\section*{Common Student Errors}
\begin{itemize}[leftmargin=*]
\item Calling a premise invalid instead of false or doubtful.
\item Calling an argument sound merely because the form is valid.
\item Saying Antony ``proves'' Caesar was not ambitious without noticing that Antony's evidence may only challenge Brutus's stronger claim.
\item Treating emotional force as if it automatically counts as evidence.
\item Forgetting that a true conclusion can be supported by bad proof.
\item Failing to identify the pressure-point premise.
\end{itemize}

\section*{Connection to Future Weeks}
Week 7 begins the fallacy sequence. Week 6 prepares it by teaching students to name the kind of failure precisely. A fallacy is not merely a false statement; it is defective reasoning that can appear to prove something. Students who can distinguish invalid form, false premise, doubtful premise, and overclaim will be ready to understand why bad arguments often feel strong.

\end{document}
